\paragraph{Boolean interval bounds.}
The cyclic bounds used above belong to a class whose greatest core
minorants and relative source compatibility have closed descriptions.
Fix distinct owners $a,b,c$, a finite nonempty passive set $V$, and a
finite family $\mathcal Z$ of intervals
\[
 \varnothing\ne L_i\subsetneq U_i\subsetneq V.
\]
The family may be empty. Write $q(L,U)$ for the atomic input with
$c$ on $L$, $a$ on $U\setminus L$, and $b$ outside $U$.
On $2^V$, draw the inclusion edges and the reversed edges $U_i\to L_i$;
write $\preceq$ for reachability. Let $\mathcal B_{\mathcal Z}$ be the
normalized monotone Boolean functions constant on these intervals.
It contains the all-variable disjunction and conjunction.
For normalized $f$, let $E_f$ be the greatest completion with
$c$-profile $f$, $b$-profile $f^\partial(A)=1-f(V\setminus A)$,
and unanimity at every other owner. Put
\[
 C_{\mathcal Z}=\sum_{f\in\mathcal B_{\mathcal Z}}E_f,
 \qquad L(w)=\{j:w_j=\{c\}\},\quad U(w)=\{j:w_j\ne\{b\}\}.
\]

\begin{theorem}[Greatest pair-erasure bound]
\label{core:thm:pair-erasure-envelope}
The operation $C_{\mathcal Z}$ is the greatest core operation satisfying
$H(q(L_i,U_i))\subseteq\{b,c\}$ for every $i$. On all winner inputs,
\begin{equation}\label{core:eq:pair-erasure-envelope}
 C_{\mathcal Z}(w)=
 \begin{cases}
 \{b,c\},&L(w)\subsetneq U(w),\ U(w)\preceq L(w),\\
 \bigcup_jw_j,&\text{otherwise}.
 \end{cases}
\end{equation}
Consequently, the operation that erases to $\{b,c\}$ precisely on the
selected interval regions and is input union elsewhere is core if and
only if every collapsed strict interval is selected.
\end{theorem}
\begin{proof}
The functions in $\mathcal B_{\mathcal Z}$ are exactly the normalized
isotone functions of $\preceq$. They all agree at $L\subseteq U$
exactly when $U\preceq L$: otherwise the reachability upset of $U$
separates them. No nonempty vertex reaches $\varnothing$, and $V$
reaches only itself, so this separator is normalized whenever needed.
The all-variable disjunction and conjunction also show that a collapsed
strict interval has nonempty lower and proper upper endpoint.

At $w$, $E_f$ returns $c$ if $f(L(w))=1$, $b$ if $f(U(w))=0$,
and input union otherwise. The only other singleton case has every
input equal to the same foreign atom and already equals input union.
A separating $f$ therefore supplies the union. On a collapsed strict
interval, disjunction supplies $c$ and conjunction supplies $b$.
If $L(w)=U(w)$, the inputs are all $b/c$ atoms and these same two
functions supply their union. This proves
\eqref{core:eq:pair-erasure-envelope} and all selected bounds.

Let $H$ satisfy the selected bounds. By the atomic-observer expansion
\eqref{core:eq:atomic-observer-expansion}, its $a$-trace is a union of
Boolean derivatives. Every component must be constant on all selected
intervals, hence on every collapsed interval. Thus $H(q(L,U))$ omits
$a$ whenever $U\preceq L$.
Replace every input of an arbitrary $w$ other than an exact $b/c$ atom
by $\{a\}$. If $H(w)$ has a winner outside $\{b,c\}$, the replacement
has winner $a$. This follows by induction: the $b/c$ singleton tests
are unchanged, so a $b/c$ root cannot acquire a new absorber; any other
root either unions its children or retains $a$ by an $a$-absorption.
Provenance excludes every other absorbing atom on the replacement.
Therefore $H(w)\subseteq\{b,c\}$ on every collapsed region, and
provenance gives $H\subseteq C_{\mathcal Z}$ elsewhere.

For the final claim, a core below the selected-region operation is
exactly a core satisfying its selected atomic bounds. Its greatest
core minorant is therefore $C_{\mathcal Z}$. A missing collapsed
interval differs at its canonical atomic row, where the selected-region
operation is $\{a,b,c\}$ and $C_{\mathcal Z}$ is $\{b,c\}$.
\end{proof}

For any core $S(X,V)$, the neutral-residue bounds give the useful equivalence
\begin{equation}\label{core:eq:pair-spine-tests}
 S\subseteq C_{\mathcal Z}+X
 \quad\Longleftrightarrow\quad
 S(\{b,c\},q(L_i,U_i))\subseteq\{b,c\}\quad\text{for every }i.
\end{equation}
Indeed $S\subseteq C_{\mathcal Z}+X$ is equivalent to
$\partial_XS\subseteq C_{\mathcal Z}$. At the selected rows the
residue's $a$-bit is that of $S(I\setminus\{a\},q_i)$, also that of
$S(\{b,c\},q_i)$ since the two ambiguous source menus have the same
$a$-view. The latter uses only $a,b,c$, so omission of $a$ is the pair
bound. This formulation works also when $K>3$.

\paragraph{Supports for the other owners.}
For $A\subseteq V$, project every selected reversed edge to
$U_i\cap A\to L_i\cap A$ on $2^A$, retaining the inclusion edges.
Define
\[
 \eta_{\mathcal Z}(A)=1
 \quad\Longleftrightarrow\quad
 A\text{ does not reach }\varnothing\text{ in this projected graph}.
\]
This is monotone in $A$, false at $\varnothing$, and true at $V$.

\begin{lemma}[Outside-owner profile image]
\label{core:lem:pair-support-profile}
For $r\notin\{b,c\}$, the profiles of core operations below
$C_{\mathcal Z}$ are exactly the normalized monotone $f$ with
$f\leq\eta_{\mathcal Z}$.
\end{lemma}
\begin{proof}
Suppose $H\subseteq C_{\mathcal Z}$ and $\Phi_r(H)(A)=1$.
In each selected row replace the free $a$ atoms by a fresh unowned
atom $\omega$. Outside-winner presence is unchanged: the $b/c$
singleton tests agree, and every other owned root acts as OR on that
presence bit, including a root that previously absorbed to $a$.
Thus the new rows also satisfy the pair bound.

Prune an $r$-winning profile proof on $A$: at an $r$-node keep a
positive child, and at every other node keep both. All remaining
leaves lie in $A$. On $r$-avoiding inputs this pruning is contractive,
as in \cref{core:lem:owner-face-erasure}. Every fresh-$\omega$ row
avoids $r$, so the resulting core on $A$ respects the projected bounds.
Its atomic-observer expansion supplies a normalized Boolean function
constant on every projected interval. Such a function exists exactly
when $A$ does not reach $\varnothing$.

Conversely, choose such a function on each feasible nonempty $A$.
Its dual-profile completion on support $A$ lies below
$C_{\mathcal Z}$ and has $r$-profile $[A\subseteq-]$.
For a prescribed $f\leq\eta_{\mathcal Z}$, fold these completions
by $\star_r$ over its positive coalitions. The fold stays below the
idempotent bound and has profile $f$.
\end{proof}

\paragraph{Four helpers replace arbitrary anchor words.}
For a full type $B$, put $\psi_r=\neg B_r$,
$\beta=\psi_c$, and $\alpha=\psi_b^\partial$.
Allow the constant functions $0,1$ in addition to normalized Boolean
functions. Let $\gamma_*$ be the reachability upset of
\[
 \{A:\alpha(A)\beta(A)=1\}
 \ \cup\ \{L_i:\alpha(L_i)=0<\alpha(U_i)=1\},
 \qquad
 \mathcal Z_0=\{i:\beta(L_i)=0,\ \alpha(U_i)=1\}.
\]

\begin{theorem}[Exact types over pair-erasure bounds]
\label{core:thm:pair-interval-types}
There is a primitive one-$X$ spine $S\subseteq C_{\mathcal Z}+X$
with full type $B$ if and only if
\[
 \gamma_*\leq\beta,
 \qquad \psi_r\leq\eta_{\mathcal Z_0}\quad(r\notin\{b,c\}).
\]
At most $K+2$ marked-path attachments suffice. For a known core
$F\in\mathcal I(C_{\mathcal Z})$, the type is actual precisely when
these conditions hold and its source box is sound for $F$.
\end{theorem}
\begin{proof}
First consider a path with owners only $b,c$.
At a canonical row, a $c$-helper with own profile $f$ can be replaced
by $E_f$, and a $b$-helper with own profile $g$ by $E_{g^\partial}$.
These replacements preserve the respective own singleton tests,
can add the other permitted singleton, and cannot add an $a$-winner.
For the last assertion, if $f(L)=0<f(U)$, the pair test excludes any
helper with that $c$-profile bounded by $\{b,c\}$ at $q(L,U)$.
Thus $E_f$'s $a$-winner was already present. The observer order on
$\{b,c,N,P\}$ is monotone for both path gates, so safety is preserved.

On passive $b/c$ inputs these canonical helpers have Boolean values.
Let $\gamma,\delta$ be the path's Boolean endpoint functions with
$X=b,c$. Induction along the path gives
\begin{equation}\label{core:eq:pair-endpoint-identities}
 \gamma=\beta\delta,\qquad \delta=\alpha\lor\gamma.
\end{equation}
At $q_i$, singleton sources are subsets of $\{b,c\}$, so safety makes
both endpoints omit $a$. The canonical helpers take the three values
$b,\{a,b,c\},c$, interpreted by Boolean endpoint pairs
$(0,0),(0,1),(1,1)$. Hence $\gamma,\delta$ are constant on every
zero interval. They satisfy $\alpha\beta\leq\gamma\leq\beta$.
The least such $\gamma$ whose $\alpha\lor\gamma$ is also constant
is exactly the indicated seed upset $\gamma_*$.

Conversely, given $\gamma_*\leq\beta$, set
$\gamma=\gamma_*$, $\delta=\alpha\lor\gamma$, and use the path
\begin{equation}\label{core:eq:four-pair-helpers}
 X\star_b E_\alpha\star_c E_\beta
   \star_b E_\delta\star_c E_\gamma.
\end{equation}
Omit a $b$-attachment with function $1$, or a $c$-attachment with
function $0$; these denote inactive factors, not constant core helpers.
All remaining functions are normalized. The total marked factors are
exact because $\alpha\land\delta=\alpha$ and
$\beta\lor\gamma=\beta$.
On a zero interval write $g\leq d$ for the common endpoint values
of $\gamma,\delta$. If $g=d=1$ the last gate resets to $c$;
if $g=d=0$ the third resets to $b$ and the last fixes it.
If $g=0,d=1$, \eqref{core:eq:pair-endpoint-identities} forces
$\alpha=1,\beta=0$ at both endpoints, so all four gates fix the
source pair. This proves safety by \eqref{core:eq:pair-spine-tests}.

Now allow all path owners. Deleting an outside-owned gate at a
canonical row can only decrease the observer value; it acts as union
on the outside bit. Thus the two-anchor condition remains necessary.
At a row in $\mathcal Z_0$ no $b/c$ helper is its own gate's atom.
Starting from the ambiguous pair, no such gate can erase an outside
winner. Therefore every outside helper respects all these pair bounds.
Its own-owner fold has profile $\psi_r$ and lies below
$C_{\mathcal Z_0}$. The preceding lemma proves necessity.

For sufficiency choose one such helper for each active outside owner
and attach them before \eqref{core:eq:four-pair-helpers}.
On $\mathcal Z_0$ the prefix stays exactly $\{b,c\}$.
At every other selected row the anchor suffix has a constant reset,
erasing the prefix. The bound is global by
\eqref{core:eq:pair-spine-tests}, and every marked factor is exact.
There are at most $K-2$ prefix attachments and four anchor attachments.
Finally use marked linearization for necessity in $F$ and soundness
with marked union completion for sufficiency, as in
\cref{core:thm:profile-interval-types}.
\end{proof}

\paragraph{The complete relative source criterion.}
The types above give an explicit interval preorder, including the
implications caused by cyclic bounds. For each strongly connected
component $M$ of the Boolean graph, retain the zero edges internal to
$M$, denoted $\mathcal Z_M$, and define the owner box
\begin{align*}
 \Theta_M(E)&=\bigwedge_r B^M_r(E_r),\\
 B^M_c(A)&=[\exists S\in M:\ A\subseteq S],\\
 B^M_b(A)&=[\exists S\in M:\ S\cap A=\varnothing],\\
 B^M_r(A)&=1-\eta_{\mathcal Z_M}(A)\quad(r\notin\{b,c\}).
\end{align*}
Its tag consists of the proper colorings
\[
 T_M=\{E:E_c,\ V\setminus E_b\in M\}.
\]
Every $T_M$ is nonempty: take a complete $b/c$ coloring corresponding
to any vertex in $M$.
The box $\Theta_M$ is zero on every other tag $T_N$: its two anchor
factors would otherwise give paths $N\preceq M\preceq N$.

\begin{theorem}[Source compatibility over pair-erasure bounds]
\label{core:thm:pair-source-interval}
A predicate $P$ on proper passive colorings is realized in
$\mathcal I(C_{\mathcal Z})$ if and only if it is antitone and,
for every component $M$,
\[
 P\text{ is positive somewhere on }T_M
 \quad\Longrightarrow\quad \Theta_M\leq P.
\]
The least realizable predicate positive at a coloring $p$ is
$\Theta_M$ if $p\in T_M$, and the principal downset of $p$ otherwise.
Each of these least predicates has a one-$X$ minorant with at most
$K+2$ marked-path attachments.
\end{theorem}
\begin{proof}
Write $C=p_c$, $D=V\setminus p_b$, so $C\subseteq D$.
For any actual type positive at $p$, its anchor endpoints satisfy
$\gamma(C)=0$ and $\delta(D)=1$.
Consequently the gap $\gamma=0,\delta=1$ contains every $A$ with
$D\preceq A\preceq C$. This set is empty if $D\npreceq C$,
and otherwise is the component $M$ containing $C,D$.

In the latter case \eqref{core:eq:pair-endpoint-identities} gives
$\beta=0,\alpha=1$ throughout $M$.
A $c$-capture on $A$ is therefore impossible when $A$ is contained
in some member of $M$; a $b$-capture is impossible when $A$ misses
some member. On internal zero edges there is no anchor reset, so the
outside profile bound gives $\psi_r\leq\eta_{\mathcal Z_M}$.
Thus every type positive at $p$ contains the entire box $\Theta_M$.
Marked linearization shows that every interval predicate positive at
$p$ must contain it.

To realize this box, put
\[
 \alpha_M(A)=[\exists S\in M:S\subseteq A],\qquad
 \beta_M(A)=[\forall S\in M:A\nsubseteq S].
\]
Choose $C_0\in M$ and set
\[
 \gamma(A)=[A\npreceq C_0],\qquad
 \delta(A)=\gamma(A)\lor[C_0\preceq A].
\]
These functions are constant on zero edges, and satisfy
$\alpha_M\beta_M\leq\gamma\leq\beta_M$ and
$\delta=\alpha_M\lor\gamma$, $\gamma=\beta_M\delta$.
Indeed below $M$ the two endpoint functions are zero, on $M$ they
are $0,1$, and outside its lower set they are both one; the definitions
of $\alpha_M,\beta_M$ give the stated identities in each case.
Moreover $\beta_M(L_i)=0,\alpha_M(U_i)=1$ holds exactly for internal
edges. The preceding type theorem, or its four-helper construction,
therefore realizes $\Theta_M$, using the outside profiles
$\eta_{\mathcal Z_M}$.
It is positive on $T_M$: for every outside color block
$A\subseteq D\setminus C$, projecting an internal path $D\to C$
gives a path $A\to\varnothing$, so its outside factor is one.
The two anchor factors are one by $C,D\in M$.

If $D\npreceq C$, take
\[
 \alpha(A)=[D\subseteq A],\quad
 \beta(A)=[A\nsubseteq C],\quad
 \gamma(A)=\delta(A)=[D\preceq A].
\]
Again the endpoint identities hold and all zero edges have equal
endpoint functions, so $\mathcal Z_0$ is empty.
Every outside profile $\psi_r(A)=[A\nsubseteq p_r]$ is then
permitted. The constructed type is exactly the principal downset box
$\bigwedge_r[A\subseteq p_r]$, positive at $p$.
Antitonicity makes this the least possible predicate.

Finally union these least boxes over the positive colorings of $P$,
and add the passive baseline. The stated conditions put every box
inside $P$, while each positive coloring belongs to its own box.
The source reconstruction formula gives exactly $P$.
\end{proof}

For example, take $|V|\geq3$ and select every nondegenerate interval.
Then $C_{\mathcal Z}(w)$ is $\{b,c\}$ whenever both a $b$-atomic and
a $c$-atomic input occur, and is input union otherwise.
All nonempty proper coalitions form one component $M$.
Every proper support collapses to $\varnothing$ in its projected graph,
whereas $V$ does not. Thus the component's activation box is simply
\[
 \Theta_M(E)=\bigwedge_r[E_r\ne V].
\]
Activating the source on any coloring with both anchor types therefore
activates it on every coloring having no monochromatic full coalition.
This baseline has unanimity singleton profiles, while its mixed-value
bounds supply the additional implication.
